\question[5]
Erkläre die drei Fehlerarten Syntaxfehler, Laufzeitfehler und logischer Fehler. 
Gib jeweils an, wann der Fehler auftritt (vor dem Start, während der Ausführung oder gar nicht als Absturz erkennbar) 
und nenne ein kurzes Beispiel.

\begin{solutionbox}{6cm}
\begin{lstlisting}
Syntaxfehler:
Das Programm verstößt gegen die Grammatik der Sprache und kann
gar nicht erst gestartet werden.
Beispiel: print("Hallo"      (Klammer fehlt)

Laufzeitfehler:
Das Programm startet, bricht aber während der Ausführung mit
einer Fehlermeldung (Traceback) ab.
Beispiel: x = 5 / 0          (ZeroDivisionError)

Logischer Fehler:
Das Programm startet, läuft bis zum Ende durch und stürzt nicht
ab, liefert aber ein falsches Ergebnis.
Beispiel: durchschnitt = 10 + 20 / 2   (statt (10 + 20) / 2)
\end{lstlisting}
\end{solutionbox}